\section{Results}
\label{sec:results}

All results come from the production design of Section~\ref{sec:stats}: the preregistered stages, with the
exact refuge capacity and the primary evaluator ($t'_{\rm pk}$, $t_{B99}$, free space in O4, O10
on adults, O11 on the second half of the births, crowded class at $m>8$; Table~\ref{tab:deviations}), plus seven
stages not in the preregistered plan; 51\,828 simulations and 960 transplant replicates. Every run was also scored with the preregistered
evaluator and with each option changed separately (Appendix~\ref{app:res-rescoring}). Unless stated otherwise, the
numbers for \textsf{C1}, \textsf{C1}$'$, \textsf{C0} and the $\alpha=0$ control come from one reference ensemble
($n=200$, $T=600$), those of \textsf{C2} from its production ensemble ($n=200$, $T=700$), and the null models, among them the long-lived
control (the core model with the four rules switched off, $\alpha=1.4$ and $a_{\max}=120$ weeks;
Section~\ref{sec:model-candidates}), start from the initial condition of \textsf{C1} (Appendix~\ref{app:res-ensembles}). $f$ denotes the
fraction of runs that pass a pattern; intervals are 95\% Wilson intervals unless another kind is named. Pass
fractions are rounded to two decimals in the tables; in the text, O17p of the $n=200$ ensembles is quoted with three
(199/200 $=0.995$, 167/200 $=0.835$, 141/200 $=0.705$), because the third decimal is exact there and separates 0.995
from 1.00.

\subsection{Production ensembles}
\label{sec:res-production}

\paragraph{\textsf{C1} reproduces every applicable essential pattern.}
Each applicable essential pattern passes in at least 99\% of the 200 runs, and the primary outcome O17p in
$f=0.995$ [0.97, 1.00]; the independent ensemble at $T=700$ gives 0.98 [0.95, 0.99], the preregistered evaluator 0.98
and the single-pass refuge admission, run as a variant, 0.995: neither the admission rule nor the
seeds change the primary outcome.
The time course (Fig.~\ref{fig:res-trajectories}) is a single boom, a plateau and an extinction by old age. The
mean social state of the adults (age $\ge6$ weeks, which includes the maturing pups of the first cohort) falls below
0.5 at week 13 [13, 14], while the population is still growing. Part of this fall is demographic dilution by pups
born at $s=0$: measured on the animals past the plasticity window ($a\ge a_w=12$), or on the founders alone, the
crossing comes at week 16, so that deterioration takes of the order of one generation, $\Gamma=t_{\rm det}/T_{\rm gen}\approx1.1$
($t_{\rm det}\approx10$ weeks, $T_{\rm gen}\approx9$; 1.2 on the founders; Section~\ref{sec:model}). The socially
functional growth phase is short either way: O10b passes in 0.02 of the runs with the adult mean of the primary
evaluator and in 0.28 (marginal in the rest) when measured on the animals past the window, a phase B of 13 weeks
that holds two thirds of the births, against 30 weeks in Universe~25. The population comes within 2\% of
its maximum at $t'_{\rm pk}=34$ [32, 36] weeks, 99\% of the births have occurred by week 40, the last surviving
birth is at week 59, the plateau lasts 51 [49, 53] weeks, and the colony then ages at one week per week and dies
out at week 169 [167, 171] (Kaplan--Meier median). The peak uses $\rho_{\rm pk}=0.36$ of the damage-threshold
capacity $K_\gamma=m_c|G|$ and leaves 16\% of the cells empty (Universe~25: 0.57 and about 20\%; an agreement in
sign only, and $f_0$ was used in the calibration); 17\% of the adults form a persistent isolated class (0.6\% in
\textsf{C0}). The desirable patterns are only partly reproduced: O8 passes in 0.62 [0.56, 0.69] of the runs, O9 in
0.45 and O13b in none (Table~\ref{tab:res-production-si}).
Without mortality before senescence the maximum $t_{\rm pk}=\arg\max N$, at week 55 [50, 60], falls on the last
birth before it in 69\% of the runs, 21 weeks after the plateau begins; every criterion normalised by the peak
time inherits this ambiguity, which is why the primary evaluator uses $t'_{\rm pk}$ and $t_{B99}$. The decline
lasts $(T_{\rm ext}-t_{\rm pk})/t_{\rm pk}=2.07$ peak times with $\arg\max N$, close to Calhoun's 1.84--2.2, but
3.97 with $t'_{\rm pk}$: this agreement is not a robust result.

\paragraph{\textsf{C1}$'$, controls and \textsf{C2}.}
\label{sec:res-c1c2}
\textsf{C1}$'$ (R1 + R2 + R7, without AV) passes O17p in 0.835 [0.78, 0.88] of the runs (0.93 with the preregistered
evaluator). What fails is O5: in 15\% of its runs the 99th percentile of the births comes when the population has
already fallen below $0.7\,N_{\rm pk}$ or later than $1.5\,t'_{\rm pk}$, because its tail of births is longer
($\tau_B=8.7$ against 6.9 weeks). Without AV the deteriorated animals gather in pools that leave half of the lattice
empty at the peak ($f_0=0.52$ against 0.16) and the senescent decline O8 passes in 0.21 instead of 0.62.
\textsf{C0} passes every instantaneous essential pattern (O17 in 0.92) but its isolated class is never persistent
(O17p $=0$): persistence is what refuges add. Without attraction ($\alpha=0$) the colony fills 71\% of $K_\gamma$
with 3\% of empty cells and fails the primary O4 and O14. The short-lived control, the core model with the four rules switched off, $\alpha=1.4$ and a short life
($a_{\max}=60$, $a_{\rm rep}=50$ weeks), fails O11 and O13 in every run, and 53\% of
its colonies survive the first collapse and oscillate; the long-lived control, identical except for $a_{\max}=120$
and $a_{\rm rep}=80$, goes extinct in
0.98 of the runs, with a peak at 0.85 of $K_\gamma$ and no failure of the late cohorts scored at six weeks of age
(scored at the end of the plasticity window they fail in 0.81 of the runs from the initial condition of \textsf{C1};
Section~\ref{sec:res-null}). Near-certain extinction is
therefore not specific to the new rules; what they add is the structure of the collapse.
\textsf{C2}, which ties the retreat to maternal socialisation, makes its isolated class a later-born cohort (O13b at
least marginal in 200/200 runs against 0/200 in the \textsf{C1} ensemble of the same stage and 2/200 in the
reference ensemble) but pays with its demography: O8 fails in every run, O5
falls to 0.815 and O4 to 0.90, and O17p is 0.705 [0.64, 0.76] against 0.98 in the independent \textsf{C1}
ensemble of the same stage ($T=700$; Holm $p=4\times10^{-14}$ over the 29 rows of the \textsf{C1}--\textsf{C2}
family, 27 distinct tests; 0.88 with the preregistered evaluator), whereas the calibration had predicted no difference
(Appendix~\ref{app:res-c1c2}).

\begin{table}[t]
\centering
\caption{Production ensembles: fraction of runs that pass each essential pattern, primary evaluator.
\textsf{C1}, \textsf{C1}$'$, \textsf{C0} and the $\alpha=0$ control: the reference ensemble ($n=200$, $T=600$);
\textsf{C2}: $n=200$, $T=700$; short-lived control ($\alpha=1.4$, $a_{\max}=60$) and long-lived control ($\alpha=1.4$, $a_{\max}=120$), both
without the four rules and from the dispersed initial condition: $n=100$, $T=1000$. Patterns are evaluated over the applicable runs; O17 and O17p over
all runs. The 95\% Wilson half-width is at most 0.07 for $n=200$ and 0.10 for $n=100$, and at most 0.02 and 0.04
for fractions of 0 or 1. No anti-pattern is triggered in the candidates; A2 (recovery of natality) is triggered in
0.02 of the \textsf{C2} runs, 0.55 of the short-lived runs and 0.03 of the long-lived runs. Magnitudes are ensemble
means, except the Kaplan--Meier median and the spread of the extinction time (interquartile range over the median
of the Kaplan--Meier estimate). Transplants: Table~\ref{tab:res-trans}.}
\label{tab:res-production}
\footnotesize
\setlength{\tabcolsep}{3pt}
\begin{tabular}{@{}llccccccc@{}}
\toprule
ID & Pattern & \textsf{C1} & \textsf{C1}$'$ & \textsf{C0} & $\alpha{=}0$ & \textsf{C2} & short-lived & long-lived \\
\midrule
O4 & peak below capacity, free space & 1.00 & 1.00 & 1.00 & 0.00 & 0.90 & 0.87 & 0.08 \\
O5 & natality ends at high $N$ & 0.99 & 0.85 & 1.00 & 1.00 & 0.81 & 0.25 & 0.90 \\
O6 & no rebound & 1.00 & 1.00 & 1.00 & 1.00 & 0.98 & 0.45 & 0.97 \\
O7 & extinction & 1.00 & 0.99 & 1.00 & 1.00 & 0.98 & 0.47 & 0.98 \\
O10 & adult deterioration before peak & 1.00 & 1.00 & 1.00 & 1.00 & 1.00 & 1.00 & 1.00 \\
O11 & late cohorts fail & 1.00 & 0.97 & 1.00 & 1.00 & 0.99 & 0.00 & 0.00 \\
O13 & two deteriorated classes & 1.00 & 1.00 & 0.92 & 1.00 & 1.00 & 0.00 & 0.00 \\
O13p & \quad persistent isolated class & 1.00 & 1.00 & 0.00 & 0.00 & 0.94 & 0.00 & 0.00 \\
O17 & Calhoun-like run (O17) & 0.99 & 0.83 & 0.92 & 0.00 & 0.72 & 0.00 & 0.00 \\
\midrule
O17p & \textbf{primary outcome} & 0.99 & 0.83 & 0.00 & 0.00 & 0.70 & 0.00 & 0.00 \\
 & \quad 95\% Wilson interval & 0.97--1.00 & 0.78--0.88 & 0.00--0.02 & 0.00--0.02 & 0.64--0.76 & 0.00--0.04 & 0.00--0.04 \\
\midrule
& $\rho_{\rm pk}=N_{\rm pk}/K_\gamma$ & 0.36 & 0.33 & 0.42 & 0.71 & 0.42 & 0.77 & 0.85 \\
& empty cells at peak, $f_0$ & 0.16 & 0.52 & 0.19 & 0.03 & 0.11 & 0.30 & 0.59 \\
& persistent isolated, $\Phi^{\rm pers}_{\rm BO}$ & 0.17 & 0.16 & 0.01 & 0.09 & 0.13 & 0.00 & 0.00 \\
& median $T_{\rm ext}$ (weeks, KM) & 169 & 182 & 166 & 155 & 218 & -- & 186 \\
& spread of $T_{\rm ext}$ (KM IQR/median) & 0.09 & 0.12 & 0.08 & 0.05 & 0.14 & -- & 0.16 \\
\bottomrule

\end{tabular}
\end{table}

\begin{figure}[t]
  \centering
  \includegraphics[width=\textwidth]{fig_res_trajectories.pdf}
  \caption{Time course of the reference ensemble ($n=200$ each).
  \textbf{(a)}~Population relative to the damage-threshold capacity $K_\gamma=m_c|G|$: median and 10--90\% band for
  \textsf{C1} (thin lines are 12 individual runs), and medians for \textsf{C1}$'$, \textsf{C0} and the $\alpha=0$
  control.
  \textbf{(b)}~Weekly births (median).
  \textbf{(c)}~Mean social state $\bar s$ of the adults (age $\ge6$ weeks; solid and dashed coloured lines as in a),
  median over living colonies, shown while at least half of the colonies are alive, on a linear scale; for
  \textsf{C1} also the mean over the animals past the plasticity window (age $\ge a_w$, dotted) and over the
  founders (dashed), which are not diluted by the newborn pups. The dotted grey lines mark the 0.5 level of O10 and
  the 0.2 deterioration threshold.
  \textbf{(d)}~Mean age, with a reference line of slope one: after natality stops, the colony simply ages.}
  \label{fig:res-trajectories}
\end{figure}

\subsection{Which patterns carry evidence: null models}
\label{sec:res-null}

A criterion is informative only if it can fail. Table~\ref{tab:res-null} gives the pass rate of every essential
pattern in four null models run from the initial condition of \textsf{C1} ($n=200$): no crowding damage
($\gamma=0$), full inheritance without social learning ($-$R2$^*$: $w=1$, individual recovery at $r=0.01$), the
long-lived demography without new rules, and \textsf{C1} without its four rules; and in the two partial ablations of
the reference ensemble ($\alpha=0$ and \textsf{C0}); Table~\ref{tab:res-null-full} (appendix) adds the desirable
patterns. NA counts as a failure. By the rule fixed before the runs, a criterion is \emph{generic} if it passes
($f\ge0.8$) in at least two of the four null models, and it \emph{discriminates} against a model if the upper
Wilson bound of its pass rate there is below 0.8 and the Newcombe interval of the difference with \textsf{C1}
excludes zero.
O5, O6, O7 and O10 are generic (they pass in three or four null models), and so is the primary O4, which still
passes in $-$R2$^*$ and without rules (the free space at the peak needs no rule: $f_0=0.74$ without rules, 0.59 in
the long-lived control; and $f_0$, part of the primary O4, was an informal target of the calibration): these five criteria
carry no evidence for the rules; they verify that the model has a boom, an end of natality and an extinction. Scored at six weeks, as in the primary evaluator, the
primary O11 discriminates against every null model, but not at the end of the plasticity window (below); O13 and
O13p discriminate against three of the four, but not
against $-$R2$^*$, where both deteriorated classes form by full inheritance without social learning (O13p 1.00
there), and O13p also discriminates against both partial ablations. O17p fails in every null model and in both
partial ablations; which criteria carry that evidence depends on the age at which O11 is scored (below). With the preregistered evaluator O10 passed in every null model
and O11 passed vacuously without crowding damage, which is why both were redefined (Table~\ref{tab:res-null-v1}).
O11 is in part definitional, because pups are born at $s=0$ and are scored at six weeks of age, inside the
plasticity window. Re-scoring the same runs with the late cohorts measured at the age $a_w=12$ instead (same
thresholds; \texttt{results\_v2/o11\_aw}), \textsf{C1} and \textsf{C1}$'$ barely change (O11 0.995 and 0.935, O17p
0.990 and 0.815), but O11 then also passes in 0.81 of the long-lived runs and in 0.755 of the runs without rules,
whose late cohorts fall below 0.2 through accumulated crowding damage, so it no longer discriminates against those
two null models; against $-$R2$^*$ it passes in 0.21 (O17p 0.205), and without crowding damage it stays NA. Their late cohorts even end lower than those of \textsf{C1} (median $\bar s$ 0.074 in the long-lived runs and
0.055 without rules, against 0.130). The dependence on $s_{\rm base}$ almost disappears at that age (below). At six
weeks, then, O11 separates \textsf{C1} from the null models largely through the birth state of the pups
($w=0.2$, $s_{\rm base}=0$; with $s_{\rm base}=0.25$ it falls to 0.08), not through a failure that the rules add to
the crowding damage, and the evidence of O17p against the long-lived
demography and the model without rules rests on O13 and O13p; against $-$R2$^*$ it rests on O11 at six weeks and is
weaker at $a_w$. O10b, the functional growth phase, passes only where crowding is absent or weak and fails in
\textsf{C1} with the adult mean (0.28 when measured on the animals past the window).

\begin{table}[t]
\centering
\caption{Discriminating power of the essential criteria (primary evaluator; fraction of all runs that pass, NA
counting as failure). Columns: the reference ensemble (\textsf{C1}, \textsf{C1}$'$, $\alpha=0$, \textsf{C0}), the
four null models with the initial condition of \textsf{C1} ($\gamma=0$; $-$R2$^*$: $w=1$ without social learning;
long-lived control; no rules: \textsf{C1} without its four rules), $n=200$ each, and the short-lived control
($n=100$). ``Generic'': passes in at least two of
the four null models. ``Discriminates'': upper Wilson bound below 0.8 and Newcombe interval of the difference with
\textsf{C1} excluding zero (long: long-lived control; none: no rules; short: short-lived control). O10b (functional growth phase) is
listed because its failure in \textsf{C1} is discussed in the text. O11 is scored at six weeks of age, as in the
primary evaluator. Scored at the end of the plasticity window ($a_w=12$; same runs, \texttt{results\_v2/o11\_aw}) it
passes in \textsf{C1} 0.995, \textsf{C1}$'$ 0.935, $\gamma{=}0$ 0.00 (NA), $\alpha{=}0$ 1.00, \textsf{C0} 0.99,
$-$R2$^*$ 0.21, long 0.81 [0.75, 0.86] and none 0.755 [0.69, 0.81] (short-lived control not re-scored): it is then not
generic (it passes in one null model) and discriminates only against $\gamma{=}0$ and $-$R2$^*$. The desirable
patterns are in Table~\ref{tab:res-null-full}.}
\label{tab:res-null}
\footnotesize
\setlength{\tabcolsep}{3.2pt}
\begin{tabular}{@{}lccccccccccl@{}}
\toprule
& \textsf{C1} & \textsf{C1}$'$ & $\gamma{=}0$ & $\alpha{=}0$ & \textsf{C0} & $-$R2$^*$ & long & none & short & generic & discriminates against \\
\midrule
O4 & 1.00 & 1.00 & 0.00 & 0.00 & 1.00 & 0.99 & 0.17 & 1.00 & 0.87 & yes & $\gamma{=}0$, $\alpha{=}0$, long \\
O5 & 0.99 & 0.85 & 1.00 & 1.00 & 1.00 & 0.97 & 0.89 & 0.91 & 0.25 & yes & short \\
O6 & 1.00 & 1.00 & 0.00 & 1.00 & 1.00 & 1.00 & 0.99 & 1.00 & 0.45 & yes & $\gamma{=}0$, short \\
O7 & 1.00 & 0.99 & 0.00 & 1.00 & 1.00 & 1.00 & 0.99 & 1.00 & 0.47 & yes & $\gamma{=}0$, short \\
O10 & 1.00 & 1.00 & 0.00 & 1.00 & 1.00 & 1.00 & 1.00 & 1.00 & 1.00 & yes & $\gamma{=}0$ \\
O11 & 1.00 & 0.97 & 0.00 & 1.00 & 1.00 & 0.00 & 0.00 & 0.00 & 0.00 & no & $\gamma{=}0$, $-$R2, long, none, short \\
O13 & 1.00 & 1.00 & 0.00 & 1.00 & 0.92 & 1.00 & 0.00 & 0.00 & 0.00 & no & $\gamma{=}0$, long, none, short \\
O13p & 1.00 & 1.00 & 0.00 & 0.00 & 0.00 & 1.00 & 0.00 & 0.00 & 0.00 & no & $\gamma{=}0$, $\alpha{=}0$, C0, long, none, short \\
O17p & 0.99 & 0.83 & 0.00 & 0.00 & 0.00 & 0.00 & 0.00 & 0.00 & 0.00 & no & $\gamma{=}0$, $\alpha{=}0$, C0, $-$R2, long, none, short \\
\midrule
O10b & 0.02 & 0.01 & 1.00 & 1.00 & 0.18 & 0.05 & 0.20 & 0.12 & 1.00 & no & none \\
\bottomrule

\end{tabular}
\end{table}


\subsection{Rule ablations and the minimal candidate}
\label{sec:res-ablations}

Fig.~\ref{fig:res-ablations} (appendix) removes the rules of \textsf{C1} singly, in pairs, in triples and all
together; each single ablation fails in a specific way. \emph{Without R7} the colony is Calhoun-like in the
instantaneous sense (O17 in 0.94), but the isolated class is never persistent (O17p $=0$). \emph{Without R2} (full
inheritance and individual recovery) the late cohorts do not fail at six weeks (O11 and O17p in 0.00 [0, 0.03]; 0.17 with the
preregistered O11, which passed when the late cohorts were empty). \emph{Without R1} adults recover when the colony
thins: natality resumes, only 23\% of the colonies go extinct by week 600, and O17p is 0.21 [0.15, 0.29].
\emph{Without AV} O17p falls to 0.79 [0.71, 0.85], through O5. Every double, triple or quadruple ablation has O17p
$=0$, and without crowding damage ($\gamma=0$) every colony grows to the cap of 15\,000 individuals. Inheritance is
not needed: with $w=0$ and a faster learning rate $r=0.15$, O17p is 1.00, so social learning is the essential part of
R2. R7 depends on its implementation and on the update scheme: without a preference ($\pi=0$) the refuges stay empty
and O17p is 0; with a soft capacity under synchronous moves O17p is also 0, because simultaneous moves send many
deteriorated animals into the same refuge, whereas under a random sequential update a soft capacity gives the same
result as the hard one (0.995 against 1.00): the hard capacity is a remedy for the collisions of the synchronous
update, not a mechanism of its own (Appendix~\ref{app:res-ablations}).

\paragraph{Q2 and the minimal candidate \textsf{C1}$'$.}
The preregistered intersection--union test is the largest of the four one-sided Fisher $p$-values, and it comes
from AV. With the primary evaluator $p_{\rm IUT}=1.1\times10^{-7}$, $p_{\rm Holm}=2.2\times10^{-7}$, so the null
hypothesis that some rule can be removed without loss is rejected; with the preregistered evaluator the
same runs give $f(-{\rm AV})=0.90$ and $p_{\rm Holm}=0.067$, not rejected, which replicates the borderline result of
the runs of the preregistered design ($p_{\rm Holm}=0.0503$). The rejection is therefore a consequence of the primary O5 alone, and specifically
of its birth reference $t_{B99}$: that option was adopted after exploratory runs with the primary evaluator had
shown that the alternative, the last sustained birth combined with $t'_{\rm pk}$, fails O5 in half of the
\textsf{C1} runs (0.48 in the production runs, Table~\ref{tab:res-rescoring}); it is a choice informed by the result
(Table~\ref{tab:deviations}), and O5 is a generic criterion (Section~\ref{sec:res-null}). The rejection of Q2 thus
shows an effect of AV on a generic criterion, significant but smaller than predicted: with either evaluator the
prediction $f(\textsf{C1}-r)\le0.5$ fails for AV (0.79), and the necessity of AV, as defined in the plan, is refuted.
The difference with \textsf{C1} is 0.20 [0.13, 0.28] in the factorial ($n=120$ per arm) and 0.16 [0.11, 0.22] in
the reference ensemble (Newcombe intervals). AV is a contributing rule: it makes the end of natality sharper (O5,
through a shorter tail of births) and supplies the senescent decline (O8 0.62 against 0.21), while it \emph{reduces}
the empty space at the peak ($f_0=0.16$ against 0.52), which is present without any rule. In \textsf{C1}$'$, R1, R2
and R7 remain necessary (O17p of 0 without each, $n=40$; Table~\ref{tab:res-c1p}).
\textsf{C1}$'$ is the minimal set \emph{conditional} on the evaluation thresholds. Re-scoring the stored class
fractions (Fig.~\ref{fig:res-theta}), O17p of \textsf{C1}$'$ is 0.83 for class thresholds $\theta\le0.125$ and 0.33
at $\theta=0.15$ (0.31 in the $-$AV arm of the factorial, $n=120$); that of \textsf{C2} falls from 0.70 to 0.01 at
0.125; \textsf{C1} holds 0.96 at 0.15. On the grid of
class and ensemble thresholds $(\theta,\tau)$, AV is ``necessary'' exactly when $\tau\ge0.8$ at $\theta\le0.125$ (0.79
lies below 0.8, with an interval that contains it) or when $\theta\ge0.15$; R1, R2 and R7 are necessary everywhere.
The necessity of R7 is necessity for O13p, a criterion designed together with R7 after \textsf{C0} had been seen;
for O11, O6 and O7 removing R7 changes nothing (O17 in 0.94).

\subsection{Confirmatory tests Q1--Q7}
\label{sec:res-Q}

Table~\ref{tab:res-Q} summarises the confirmatory family with the primary evaluator; Table~\ref{tab:res-Qcmp}
repeats it with the preregistered evaluator. With the primary evaluator six
of the seven null hypotheses are rejected after Holm adjustment and the quantitative predictions are met for Q1, Q3,
Q4 and Q7 and not for Q2, Q5 and Q6; with the preregistered evaluator five are rejected and four predictions hold
(Q3, Q4, Q6 and Q7). The two counts differ in three places: Q2 is rejected only with the primary O5 (below), the
prediction of Q1 holds only at the earlier transplant time $1.75\,t'_{\rm pk}$, and the prediction of Q6 holds only
with the preregistered O11. Q5 is refuted in the opposite direction with both. Of the predictions that hold, Q1 is a
consistency check of R1 and Q4 of the demography, and the contrast of Q3 with the short-lived control is confounded by the
demography. A rejected null hypothesis is not the same as a Calhoun-like outcome, as Q7 shows.

\begin{table}[t]
  \centering
  \caption{Preregistered confirmatory family with the primary evaluator (hypotheses and tests in
  Section~\ref{sec:stats}; Holm correction at $\alpha=0.05$ over seven tests; the $p$ of Q6 is the quasi-binomial
  one). ``Rejected'' refers to the null hypothesis after Holm adjustment; ``Met'' to the quantitative prediction
  fixed before the runs. $X$ is the cross-transplant arm. The same family with the preregistered evaluator is in
  Table~\ref{tab:res-Qcmp}.}
  \label{tab:res-Q}
  \footnotesize
  \setlength{\tabcolsep}{4pt}
  \begin{tabular}{@{}lp{4.0cm}p{5.4cm}llcc@{}}
    \toprule
    & Prediction & Estimate & $p$ & $p_{\rm Holm}$ & Rejected & Met \\
    \midrule
    Q1 & $P_T(\mathrm{D}\mid\textsf{C1})\le0.05$, $P_T(\mathrm{D}\mid X)\ge0.25$ & 0/60 vs.\ 51/60 ($P_T=0.85$ [0.74, 0.92]); discordant pairs 51:0 & $4.4\times10^{-16}$ & $1.3\times10^{-15}$ & yes & yes \\
Q2 & $f(\textsf{C1})\ge0.9$, $f(\textsf{C1}-r)\le0.5\ \forall r$ & 0.99 vs.\ $-$R1 0.21, $-$R2 0.00, $-$AV 0.79, $-$R7 0.00 & $1.1\times10^{-7}$ & $2.2\times10^{-7}$ & yes & no \\
Q3 & $P_{\rm ext}\ge0.99$, CV $\le0.10$; short-lived $\le0.7$ & 200/200 vs.\ 47/100; CV 0.066 [0.059, 0.071] & $<10^{-16}$ & $<10^{-16}$ & yes & yes \\
Q4 & slope $\in[0.8,1.25]$, $|\bar r|\le10$ wk & slope $0.996\pm0.014$; $\bar r=-2.0\pm0.3$ wk & $<10^{-16}$ & $<10^{-16}$ & yes & yes \\
Q5 & $b_2<0$, CI of $n^*_{\rm opt}\subset[4,12]$ & $b_2=+185.3\pm64.3$ (convex); no interior maximum & 1.00 & 1.00 & no & no \\
Q6 & $0<\theta<1$ & $\hat\theta=-0.15$ [-0.20, -0.05] (quasi-binomial, $\phi=3.1$) & $<10^{-16}$ & $<10^{-16}$ & yes & no \\
Q7 & corner $\ge0.8$, \textsf{C1} $\le0.2$ & 120/120 vs.\ 0/30 (30/30 explode) & $<10^{-16}$ & $<10^{-16}$ & yes & yes \\
\bottomrule

  \end{tabular}
\end{table}

\paragraph{Q1 and the transplants.}
Four males and four females of each \textsf{C1} colony, taken at $1.75\,t'_{\rm pk}$ (week 60; ages 37--64), never
found a colony in a \textsf{C1} enclosure (0/60), but do in 51/60 trials in an enclosure without R1 and R2 with
$r=0.1$ (arm $X$; exact McNemar $p=4\times10^{-16}$); the prediction ($\ge0.25$ in $X$) holds. The same animals have
$s\approx0$ (median 0.000, maximum 0.36) and fecundity scales as $s_f^2s_m$, so under R1 they cannot found a colony
with any partner (0/120 with healthy partners): a consistency check of R1, not a finding about the mechanism. The
rate in arm $X$ is a matter of age: at $1.75\,t_{\rm pk}=$ week 97 (preregistered protocol; ages 59--80) the same
colonies give 12/60, and at week 95 in every arm, the common time that removes the confound between rule and
sampling time, 5/60 (Appendix~\ref{sec:res-transplant}). An additional arm samples adults past the plasticity window with
$0.2\le s\le0.5$ (median 0.36, weeks 8--11): under R1 they found a colony in 5/60 trials, against 57/60 for the same
animals in arm $X$ and 50/60 without R1, and in the colonies they do found the adult descendants settle at the
social state of their parents (0.28--0.45 in \textsf{C1}, 0.33--0.37 in $X$). By the rule fixed in advance they
recover partly: the irreversibility of R1 + R2 is strict only for $s\approx0$, and for intermediate animals R1
acts through fecundity, not through the lineage.

\paragraph{Q3 and Q4.}
\textsf{C1} went extinct in 200 of 200 runs by $T=600$, against 47 of 100 runs of the short-lived control by $T=1000$ (one-sided
Fisher $p<10^{-16}$), with a coefficient of variation of $T_{\rm ext}$ of 0.066 [0.059, 0.071], as predicted. The
contrast that isolates the rules is the long-lived demography alone, which goes extinct in 0.98--1.00 of the runs
($p=0.11$): extinction does not discriminate. The Kaplan--Meier interquartile range over the median of
$T_{\rm ext}$ is 0.089 [0.071, 0.105] in \textsf{C1}, 0.160 [0.144, 0.180] in the long-lived control with the same
initial condition and 0.121 [0.106, 0.148] in \textsf{C1}$'$, whose difference from the control, $-0.040$
[$-0.065$, $-0.006$] (percentile bootstrap over runs), is \emph{not} within the preregistered equivalence margin of $\pm0.05$, which the runs of the preregistered design had met. A narrow
spread is not a signature of the rules (0.045 without attraction, 0.084 in \textsf{C0}): it is set by how abruptly
natality stops and by the lattice size (below). The time from the last birth to extinction grows one-to-one with
$a_{\max}$ (slope $0.996\pm0.014$ over 300 runs, TOST $p<10^{-16}$), as the demography requires
(Eq.~\eqref{eq:phaseD}; Appendix~\ref{app:res-Q}).

\paragraph{Q5 and Q6.}
The coexistence band predicted at $n^*=1/\kappa-1/\alpha\approx m_c$ (Eq.~\eqref{eq:model-nstar}) is refuted:
O13 passes in at least 95\% of the runs at 46 of the 49 points of the $(\alpha,\kappa)$ plane with $n^*>0.5$, and
the fitted quadratic in $\log n^*$ is convex; with refuges the isolated class no longer depends on the balance
between attraction and aversion. The prediction of Q6, $0<\theta<1$ for the Calhoun-like boundary in
$\log r+\theta\log a_w$ on the plane $(r,a_w)$, is not supported with the primary evaluator. The boundary is nearly
vertical in $r$: O17p is at least 0.67 for $r\le0.10$ and at most 0.60 for $r\ge0.2$ at every window, so the window
has little or no effect, and the single-index fit gives $\hat\theta=-0.15$ [$-0.20$, $-0.05$] (profile likelihood,
quasi-binomial correction $\phi=3.1$) in a model that an additive one beats ($\Delta{\rm QAIC}=28$); with the
preregistered evaluator the single-index model is the preferred one ($\Delta{\rm QAIC}=11$ in its favour) and gives
$\hat\theta=0.30$ [0.15, 0.50], as in the runs of the preregistered design (0.35). The small negative value should not be read as a reversal of
the predicted dependence: it comes from the primary O11, which at fast learning ($r\ge0.2$) fails more often the
shorter the window (O11 passes in 0.10 of the runs at $r=0.2$ with $a_w=4$ and in 0.87 with $a_w=24$), and the
change of criterion also changes which model fits the plane. The elasticity argument of
Section~\ref{sec:meanfield} does not describe this plane.

\paragraph{Q7: endogenous collapse without crowding, and the dependence on $\Lambda=r\,a_w$.}
With crowding damage switched off ($\gamma=0$), the four subcritical points ($w\le0.1$, $r\le0.03$) went extinct in
120/120 runs, while all 30 \textsf{C1} colonies exploded (one-sided Fisher $p<10^{-16}$); the colonies of the corner
first overshoot the damage-threshold capacity ($\rho_{\rm pk}=0.91$--$1.60$), so none is Calhoun-like. With
$a_w\in\{6,12,24\}$ varied at $\gamma=0$ (54 points, 20 runs each; Fig.~\ref{fig:res-q7aw},
Table~\ref{tab:res-q7aw}), 51 of the 54 points have $P_{\rm ext}=0$ or 1, so the logistic model in
$\log r+\theta\log a_w$ fixed before the production runs is separated: its profile intervals and likelihood ratio measure the
geometry of the grid, not sampling error, and the planned equivalence test of the two exponents (TOST, margin
$[2/3,3/2]$) cannot be carried out as inference (Appendix~\ref{app:res-Q}). What the design does establish is
non-parametric. The transition in $\Lambda=r\,a_w$ is bracketed by two adjacent grid points, a factor 1.25--1.5
apart; at every inheritance weight the curves for $a_w=12$ and 24 lie in the same bracket and the curve for $a_w=6$
one grid step below ($\Lambda_{1/2}\in(0.60,0.84)$ against $(0.84,1.2)$ at $w=0$; log-linear interpolation 0.71,
0.99 and 1.00). The exponent implied between windows is therefore in $(0,1)$ between 6 and 12 weeks (interpolated
0.5--0.7 over the three $w$) and compatible with 1 between 12 and 24 weeks (interpolated 0.88--1.00), where the grid
bounds it only within about $(0.5,1.5)$. The firm result is qualitative: at equal product $r\,a_w$ the six-week
window is less effective, so a single exponent does not describe the three windows and the effective learning grows
less than linearly with the window at short windows. The boundary fitted as $\sigma w+\varphi\Lambda=1$ gives
$\hat\varphi\approx1.1$--1.2 in both planes, above the preregistered $\varphi\le1$, but it is an extrapolation of the
same separated predictor and carries no valid interval.

\subsection{Robustness}
\label{sec:res-robustness}

\paragraph{Demographic and update variants, and the initial condition.}
Fig.~\ref{fig:res-rob} and Table~\ref{tab:res-rob} report \textsf{C1} and \textsf{C1}$'$ under the robustness variants
($n=200$ each), scored with both evaluators and also with O17p restricted to the criteria
that discriminate against the null models (O11 scored at six weeks, O13, O13p and the anti-patterns); by the rule fixed before the runs, a
variant preserves the result if O17p $\ge0.8$. With a constant background hazard $\mu_{\rm bg}$ at all ages the
plateau shrinks (51 weeks at 0, 35 at 0.001, 17 at 0.003 and 10 at 0.008 per week) and O17p falls to 0.83, 0.63,
0.55 [0.48, 0.62] and 0.06 in \textsf{C1} (0.67, 0.38, 0.28 and 0.04 in \textsf{C1}$'$; 0.45 at 0.003 with the
preregistered evaluator). The only criterion that fails is O5, and it fails through its timing clause, not through its
level: at $\mu_{\rm bg}=0.003$ O5 is marginal in 83 of the 200 runs and fails in 7, because the start of the plateau
$t'_{\rm pk}$ moves from week 34 to 29 while $t_{B99}$ moves from 40 to 43, so $(t_{B99}-t'_{\rm pk})/t'_{\rm pk}$
crosses the 0.5 threshold of O5 while natality still ends with the population near its maximum; at
$\mu_{\rm bg}=0.008$ O5 fails outright in 70 runs (marginal in 118). The classes, O11 and the low peak are
unchanged, and the restricted outcome is 1.00 at every $\mu_{\rm bg}$. Since O5 is generic (Section~\ref{sec:res-null}),
what the background mortality removes is the generic part of the primary outcome, not the part that carries
evidence for the rules; the primary outcome as defined nevertheless holds only without, or with very weak,
background mortality, which is stated as a limitation (Section~\ref{sec:limitations}). With initial ages drawn
uniformly from 4--52 weeks instead of all equal, O17p is 0.93 [0.89, 0.96]; from 4--80 weeks, which includes
post-reproductive founders, 0.63 [0.56, 0.69] (0.68 and 0.31 in \textsf{C1}$'$), again through the timing clause of
O5 (73 marginal and 1 failed of 200; plateau 28 and 22 weeks, $t'_{\rm pk}$ at weeks 30 and 29) and with the
restricted outcome at 1.00. The timing of the end of natality relative to the start of the plateau is therefore a
property of the synchronous initial condition and of the absence of background mortality, not of the rules; the end
of natality near the maximum itself persists. The other variants preserve the result: a nest period of three weeks
gives 0.835 [0.78, 0.88], borderline because the interval contains 0.8 (0.60 in \textsf{C1}$'$; O11 0.90,
$\rho_{\rm pk}=0.31$); a random sequential update gives 1.00
with the hard and 0.995 with the soft capacity, with $\rho_{\rm pk}=0.43$; and the exact
refuge capacity gives 0.995 against 0.995 with the single-pass admission. Pups born above the
deterioration threshold ($s_{\rm base}=0.25$ or 0.5 instead of 0) give O17p of 0.08 and 0.00 (0.64 and 0.36 with
the preregistered O11), entirely through O11 scored at six weeks of age: the failure of the late cohorts at that age is
partly definitional, tied to $s_{\rm base}=0$. Scored at the age $a_w$, the same runs give O11 of 0.975 and 0.785
and O17p of 0.965 and 0.76 (0.755 and 0.77 in \textsf{C1}$'$): by the end of the window the late cohorts are below
0.2 whatever their birth state, through crowding damage and learning from deteriorated neighbours.

\paragraph{Lattice size, seeding density and phase planes.}
Over $L=15$--60 at fixed density $N_0/L^2=0.44$ ($n=60$; Table~\ref{tab:res-size}), the peak density (0.35--0.37)
and the persistent fraction (0.17--0.18) do not change, the empty fraction decreases slightly (0.18 to 0.15), and
O17p rises from 0.82 [0.70, 0.89] at $L=15$ to 1.00 at $L\ge45$; the losses at $L\le20$ come from O5 (5 of 60 runs
at each size: a relatively longer tail of births with $N_0=100$--178), from O11 (1 of 60 at each) and, at $L=15$
only, from the anti-pattern A2, a recovery of natality in 6 of 60 runs, and the interval at $L=15$ contains 0.8.
The spread of the extinction time falls from 0.15
at $L=15$ to 0.08 at $L\ge30$ together with the relative spread of the last birth (0.58 to 0.18), while
$T_{\rm ext}-t_{\rm last}$ stays at 110--112 weeks: the spread of $T_{\rm ext}$ is that of the end of natality, as
Eq.~\eqref{eq:phaseD} implies, and its value is a property of the system size. Over seeding densities from 0.02 to
3.6 animals per cell the peak density is monotone in $N_0/L^2$ and O17p reaches 0.8 only for
$0.2\lesssim N_0/L^2\lesssim0.9$; with Calhoun's four founding pairs no run is Calhoun-like, and the best point of
the founder plane reproduces O1--O3 with O17p in 0.43 [0.27, 0.61] of the runs (Appendix~\ref{app:res-p4p6}).
Around \textsf{C1} the primary outcome holds over part of the explored planes (Appendix~\ref{sec:res-planes}): with
crowding damage every point of the $(r,w)$ plane goes extinct and O17p is at least 0.83 for $w\le0.3$ and
$r\le0.10$, falling at $w\ge0.5$ and $r\ge0.2$ through O11; every finite window $a_w\le24$ gives certain
extinction; in the $(\alpha,\kappa)$ plane O17p is at least 0.8 at 34 of the 48 points with $\alpha\ge2$, failing
where the free space of the primary O4 is lost; and a persistent isolated class requires a strong preference and a
total refuge capacity $f_Rc_R\ge0.4$. These are point estimates with 30 runs per cell (Wilson half-width up to
$\pm0.17$); re-running the 40 points nearest to a boundary with 100 runs changed them by 0.07 on average.

\paragraph{Global: Latin hypercubes and variance decomposition.}
Over two independent 15-factor Latin hypercubes (400 points, 10 runs each, plus 40 runs at the 18 points near the
boundary), 2.0\% [0.7, 3.6] of the points are Calhoun-like with the primary evaluator ($k/n\ge0.8$; bootstrap
interval over points); the mean of O17p over the box, an unbiased estimate of the volume weighted by the
probability of being Calhoun-like, is 4.0\% [2.6, 5.9], and the empirical-Bayes fraction of points with $p\ge0.8$,
which shrinks each point towards a logistic emulator, is 1.8\% [0.7, 3.2], with 8 points still uncertain
(Table~\ref{tab:res-region}, Fig.~\ref{fig:res-region}). With the preregistered evaluator the same runs give 7.2\%,
11.0\% and 6.9\% [4.8, 9.2], the values of the runs of the preregistered design; the drop comes mostly from O5 with $t'_{\rm pk}$ and from the
free-space condition of O4. Every such fraction is relative to the box and to the sampling scales of
Table~\ref{tab:params}. Part of the small fraction is definitional and part is mechanism: with $c_R=3$ refuge
residents exceed the isolation threshold and none of the 100 such points is Calhoun-like; restricting successively
to $c_R\le2$, $m_c\le10$, $f_Rc_R\ge0.3$ and $\pi\ge100$ (chosen after the runs) raises the fraction to 2.7\%, 3.1\%,
5.6\% and 10.8\% [2.7, 21.6] (4 of 37 points; 46\% with the preregistered evaluator), and even then 33 of the 37 points
fail, through O11 (16), the persistence of the class (14), O5 (9) and O4 (9); with class thresholds of 5\%, 10\% and
15\% the fraction is 5.5\%, 2.3\% and 0.5\%. The Sobol' decomposition uses 11 factors (the 8 selected by the Morris
screening plus $w$, $r$ and $a_w$; Saltelli design, $N=256$, 6 runs per point), so its indices are conditional on
four factors held at their \textsf{C1} values ($f_R$, $\pi$, $a_{\max}$, $\delta$). O17p is zero at 78\% of its
points; its total indices are led by $c_R$ ($0.49\pm0.16$), $w$ (0.43), $\kappa$ (0.41), $r$ (0.32) and $m_c$ (0.31),
which cannot be ordered, and its first-order indices add up to only 0.15, so it is almost entirely interactions;
the continuous outputs are nearly additive ($\rho_{\rm pk}$: $r$, $\kappa$, $\eta$, $\gamma$; $f_0$: $\kappa$, $m_c$;
$\Phi^{\rm pers}_{\rm BO}$: $c_R$, $m_c$). Appendix~\ref{sec:res-sensitivity} gives the full decomposition.
